% sec_technologies.tex — Vitrine v5: Technologies and Related Work
% Capture-adaptive video→3D→UE 5.8 pipeline.
% Included by the v5 master document.
%
% British English throughout.
% All \cite{} keys are defined in report/v5/references.bib.

\section{Technologies and Related Work}
\label{sec:tech}

Vitrine assembles a heterogeneous stack that spans classical Structure-from-Motion
(SfM), neural radiance fields, volumetric 3D Gaussian representations, generative
diffusion models, and real-time game-engine delivery.  This section surveys each
component in pipeline order, placing Vitrine's choices in the context of the
relevant literature and explaining the motivation for each decision.

% ─────────────────────────────────────────────────────────────────────────────
\subsection{Scene Reconstruction: from Radiance Fields to Gaussians}
\label{sec:tech:recon}

\paragraph{Neural Radiance Fields.}
The modern era of photogrammetric reconstruction from uncontrolled video begins with
\emph{NeRF} \cite{mildenhall2020_nerf}, which encodes a scene as a continuous
volumetric function and synthesises novel views by differentiably integrating colour
and density along camera rays.  NeRF established the benchmark for view-synthesis
fidelity but requires minutes-to-hours of training per scene and cannot render in
real time.  Mip-NeRF~360 \cite{barron2022_mipnerf360} extended the approach to
unbounded outdoor scenes with anti-aliasing, but the fundamental throughput barrier
remained.

\paragraph{3D Gaussian Splatting.}
Kerbl~et~al.\ \cite{kerbl2023_3dgs} replaced the implicit MLP with an explicit set
of anisotropic 3D Gaussians, each carrying position, covariance, opacity, and
spherical-harmonic colour coefficients.  Differentiable tile-based rasterisation
renders the scene at real-time frame rates while achieving view-synthesis quality
competitive with the best NeRF variants on standard benchmarks.  The explicit,
editable representation and the absence of a ray-marching inner loop make 3DGS the
natural foundation for a practical pipeline.  Vitrine uses the
\emph{LichtFeld~Studio} \cite{lichtfeld2026} implementation, which adds
ImprovedGS\texttt{+} densification heuristics and a native C\texttt{++}23/CUDA
training loop exposed through a 70-tool MCP server.

A key limitation of 3DGS for game-engine delivery is its incompatibility with
standard rasterisation pipelines: Gaussians cannot be imported directly as
collision-ready meshes or LOD assets.  NanoGS \cite{nanogs2026}, released in March
2026, demonstrated a Nanite-inspired plugin for Unreal Engine~5 that applies
hierarchical LOD clustering and GPU-accelerated depth sorting to render splat point
clouds inside UE5 at interactive rates.  Vitrine uses NanoGS as an optional
``splat-in-UE'' path for scene preview, while polygonal meshes remain the primary
game-asset deliverable.

\paragraph{Structure from Motion and Feature Matching.}
Camera pose estimation is the prerequisite for all downstream reconstruction.
Vitrine uses COLMAP~\cite{schoenberger2016sfm, schoenberger2016mvs}, the de-facto
open-source SfM framework, for sparse reconstruction and multi-view stereo.  COLMAP
was historically paired with SIFT~\cite{lowe2004_sift} features; Vitrine replaces
SIFT with \emph{ALIKED}~\cite{zhao2023_aliked}, a lighter keypoint and descriptor
extractor built on a sparse deformable descriptor head that matches dense features
at a fraction of the memory cost.  Correspondence is established with
\emph{LightGlue}~\cite{lindenberger2023_lightglue}, an adaptive Graph Neural
Network matcher that prunes the attention graph based on estimated match difficulty,
achieving SuperGlue-level accuracy at substantially lower latency.  Together,
ALIKED+LightGlue represent the current state of the art in efficient
sparse matching and have been shown to improve downstream 3DGS quality on
challenging indoor captures.

% ─────────────────────────────────────────────────────────────────────────────
\subsection{Surface Meshing}
\label{sec:tech:meshing}

Gaussian splatting does not produce a mesh natively.  Vitrine requires polygonal
geometry for Unreal Engine collision, Nanite virtualisation, and UV-baked texture
export.  Several approaches have been proposed to bridge the gap.

\paragraph{Geometry-regularised Gaussian variants.}
\emph{SuGaR}~\cite{guedon2024_sugar} (Surface-aligned Gaussian splatting) adds a
regularisation term that encourages Gaussians to align with the scene surface,
enabling Poisson mesh extraction.  \emph{2DGS}~\cite{huang2024_2dgs} collapses 3D
Gaussians onto 2D oriented planar disks, recovering surface normals analytically and
facilitating TSDF-style marching-cubes extraction.  \emph{PGSR}~\cite{chen2024_pgsr}
(Planar-based Gaussian Splatting) pursues a similar idea in the TVCG setting,
representing local geometry via planar primitives regularised for planarity
consistency.  \emph{GOF}~\cite{yu2024_gof} (Gaussian Opacity Fields) derives a
volumetric opacity field from a ray-tracing formulation of Gaussian rendering,
enabling direct level-set mesh extraction via marching tetrahedra without Poisson
post-processing.

\paragraph{Post-hoc mesh extraction.}
\emph{MILo}~\cite{guedon2025_milo} (Mesh-In-the-Loop Gaussian Splatting) takes a
fundamentally different approach: it extracts a differentiable mesh at \emph{every
training iteration} and back-propagates mesh gradients into the Gaussian parameters,
enforcing bidirectional consistency between the volumetric and surface
representations.  Published as a SIGGRAPH Asia 2025 journal track paper (ACM TOG),
MILo produces the finest geometric detail of any extraction method tested on
Vitrine's captures and is used as the \emph{primary object-mesh backend}.

\emph{CoMe}~\cite{radl2026_come} (Confidence-based Mesh extraction) assigns a
learnable confidence value to each Gaussian, using it to balance photometric and
geometric supervision.  CoMe avoids the need for multi-view normal networks or
iterative extraction and runs substantially faster than MILo.  It serves as
Vitrine's \emph{primary scene-environment mesh backend}, with a TSDF fallback
(below) for scenes where CoMe confidence fails to converge.

\paragraph{TSDF fallback.}
Truncated Signed Distance Function fusion~\cite{curless1996_tsdf} with marching
cubes~\cite{lorensen1987_marching} integrated via \emph{Open3D}~\cite{zhou2018_open3d}
provides a geometry-agnostic fallback.  The \texttt{ScalableTSDFVolume} pipeline
integrates depth images rendered from the trained splat, producing watertight meshes
even for wide-baseline captures where Gaussian regularisation fails.  Vitrine uses
the legacy \texttt{ScalableTSDFVolume} path (rather than the newer VoxelBlockGrid)
for stability on large scenes.

% ─────────────────────────────────────────────────────────────────────────────
\subsection{Per-Object Reconstruction}
\label{sec:tech:objects}

Reconstructing individual objects to game-asset quality from a room-scale scan
requires segmentation, generalised-prior 3D generation, and texture synthesis.

\paragraph{Instance segmentation.}
Vitrine uses the \emph{Segment Anything Model}~(SAM)~\cite{kirillov2023_sam} and
its video-capable successor \emph{SAM~2}~\cite{ravi2024_sam2} for per-frame mask
propagation across the capture sequence.  SAM's promptable, class-agnostic
segmentation enables instance crops without domain-specific training.  The
\emph{SAM-3D~Objects} library~\cite{sam3dobjects2025} lifts SAM masks into 3D using
the reconstructed point cloud, providing per-object bounding volumes and concept
labels that drive downstream hull reconstruction.

\paragraph{Image-to-3D generation (primary path).}
\emph{TRELLIS.2}~\cite{xiang2025_trellis2} is the designated primary hull backend.
Developed at Tsinghua University and Microsoft Research and released under the MIT
Licence, TRELLIS.2 introduces O-Voxel (Omni-Voxel), a field-free sparse voxel
representation that encodes both geometry and physically-based rendering (PBR)
material attributes.  A 4B-parameter flow-matching model trained on diverse 3D asset
datasets generates fully textured meshes (base colour, metallic, roughness, opacity)
in 17--60\,s per object on an RTX~6000~Ada, and accepts up to 16 multi-view inputs,
directly consuming the image crops Vitrine provides.  It builds on the original
\emph{TRELLIS}~\cite{xiang2024_trellis} (CVPR~2025 Spotlight), which introduced the
SLAT (Structured Latent) representation for unified Gaussian/mesh/radiance-field
decoding.

\paragraph{Image-to-3D generation (secondary / hero-asset path).}
\emph{Hunyuan3D~2.0}~\cite{zhao2025_hunyuan3d20} and
\emph{Hunyuan3D~2.1}~\cite{hunyuan2025_21} from Tencent provide the highest
open-weight texture fidelity (up to 8K PBR), at the cost of higher latency
($\approx$139\,s per object).  The system comprises a shape-generation diffusion
transformer (Hunyuan3D-DiT) and a texture-synthesis model (Hunyuan3D-Paint),
separately invocable as ComfyUI nodes.  Vitrine uses Hunyuan3D-2.1 as the
secondary/hero-asset backend, activating it when an object is designated a focal
point of the scene.

\paragraph{Inpainting and view completion.}
Real captures invariably include occluded object faces that neither TRELLIS.2 nor
Hunyuan3D can hallucinate from the available views.  Vitrine uses two complementary
models.  \emph{FLUX.2-dev}~\cite{flux2025} (Black Forest Labs) is a 24B-parameter
rectified-flow transformer with a Mistral-3 text encoder, providing the strongest
available long-form prompt fidelity for rich appearance description.
\emph{Qwen-Image-Edit}~\cite{qwenvl2025} (Alibaba) is an instruction-following
visual-language model that uniquely supports instruction-driven view rotation, making
it the preferred recovery model for occluded-face completion and the only
commercially-licensed option in the stack.

\paragraph{View-generation conditioning.}
\emph{MV-Adapter}~\cite{mvadapter2024} (Apache-2.0) generates camera-conditioned
orbit sheets via ControlNet-compatible 6-view SDXL conditioning.  It provides a
lightweight, commercially-licensed alternative to non-commercial orbital-video
models for explicit camera-parametrised multi-view synthesis.

% ─────────────────────────────────────────────────────────────────────────────
\subsection{Reconstruction Enhancement: ArtiFixer}
\label{sec:tech:artifixer}

Even a dense splat contains structurally unobserved regions --- surfaces the camera
never approached --- where 3DGS places floaters, ghosts, or holes.  These defects
propagate into the downstream mesh and UE asset.
\emph{ArtiFixer}~\cite{delutio2026_artifixer}, from the NVIDIA Spatial Intelligence
Lab and presented at SIGGRAPH~2026, is an auto-regressive video-diffusion refiner
built on \emph{Wan2.1-T2V-14B}~\cite{wan2025}.  Its opacity-mixing trick
($z_{\mathrm{mix}} = O_z \cdot z_{\mathrm{deg}} + (1-O_z) \cdot \varepsilon$)
leaves observed pixels ($O_z \approx 1$) faithful to the input, while generating
freely in unobserved regions ($O_z \approx 0$) using a strong video-diffusion prior.
The \emph{ArtiFixer3D} variant distils the generated views back into an explicit
3DGRUT~\cite{wu2025_3dgut} Gaussian recon --- a cleaner splat whose downstream mesh
is measurably less noisy than the raw LichtFeld output on coverage-limited captures.

Vitrine forks ArtiFixer to run on RTX~6000~Ada (\texttt{sm\_89}) by replacing the
FlashAttention-3/4 build block with cuDNN SDPA, with no application-code changes
(ADR-020; see \S\ref{sec:method:overview}).  The enhanced recon branch is
\emph{optional and per-scene gated}: it activates only when a quality classifier
predicts coverage-limited rather than blur-limited failure, because ArtiFixer does
not deblur --- it stays faithful ($O_z \approx 1$) in observed-but-blurry regions.

% ─────────────────────────────────────────────────────────────────────────────
\subsection{Capture Quality and Frame Selection}
\label{sec:tech:quality}

The dominant quality bottleneck in real-world captures is motion blur, not
reconstruction arithmetic.  Vitrine addresses this at the frame-selection stage
before any 3D optimisation begins.

\paragraph{No-reference image quality.}
\emph{MUSIQ}~\cite{ke2021_musiq} (Multi-scale Image Quality Transformer) is a
patch-based transformer that processes full-size images at their native resolution,
using a hash-based 2D spatial embedding and scale embedding to capture image quality
at multiple granularities without fixed-size preprocessing distortion.  Trained on
PaQ-2-PiQ, KonIQ-10k, and SPAQ, MUSIQ achieves state-of-the-art no-reference image
quality assessment (NR-IQA) and is used as Vitrine's per-frame quality gate.
\emph{Q-Align}~\cite{wang2023_qalign} provides an LMM-based alternative that
emulates discrete human rating levels; it is evaluated as a higher-cost option
for batch scoring.  Both are accessed through \emph{pyiqa}~\cite{chen2022_pyiqa},
an open-source PyTorch IQA toolbox.

\paragraph{Focus and sharpness.}
Low-level sharpness is measured with the variance-of-Laplacian operator, a classical
focus measure that computes the variance of the Laplacian-filtered image; high
variance corresponds to sharp edges.  Pertuz~et~al.\ \cite{pertuz2013_focus}
provide a comprehensive survey and comparison of such operators.  The
full-resolution variant (applied on the full 4K frame rather than a downsampled
proxy) is used, as Vitrine's capture experiments showed that downsampling to 960p
compresses the dynamic range of the sharpness score by a factor of~105, masking
meaningful differences between frames.

\paragraph{Deblur-aware Gaussian training.}
When the frame-quality gate accepts footage that is mildly blurry rather than
sharply filtered, \emph{BAD-Gaussians}~\cite{zhao2024_badgaussians} (Bundle
Adjusted Deblur Gaussian Splatting, ECCV~2024) provides a training-time deblur
path.  BAD-Gaussians models the camera trajectory during exposure as a sequence of
virtual sharp cameras, blending their renders to simulate blur formation, and
jointly optimises the trajectory and the Gaussian parameters.  This yields both a
deblurred splat and corrected camera poses.

% ─────────────────────────────────────────────────────────────────────────────
\subsection{Texture Baking and Mesh Post-processing}
\label{sec:tech:texture}

Once a mesh is available from CoMe, MILo, or TSDF, Vitrine prepares it for
game-engine import.

\paragraph{UV atlas generation.}
Vitrine uses \emph{xatlas}~\cite{xatlas2018} for automatic UV atlas generation.
xatlas is a lightweight C\texttt{++}11 mesh-parameterisation library derived from
Thekla Atlas that segments meshes into charts, parameterises each chart, and packs
them into one or more non-overlapping atlases suitable for texture baking and
lightmap generation.

\paragraph{Texture baking.}
Vertex colours from the trained Gaussian splat are baked into UV-parametrised
textures using Blender's~\cite{blender2024} Cycles GPU renderer via the
\texttt{texture\_baker.bake\_from\_vertex\_colors} pipeline with \texttt{smooth=0}
(MeshCleaner smoothing disabled to preserve sharp architectural edges).  The output
is an FBX file with baked albedo textures, imported into Unreal Engine~5.8 as a
Nanite-enabled game asset.

\paragraph{Mesh cleaner.}
Raw CoMe and TSDF meshes contain stray bridging triangles, sub-threshold connected
components, and degenerate faces.  The \texttt{mesh\_cleaner} module applies
edge-length filtering followed by connected-component analysis (retaining only the
largest $k$ components) to reduce the mesh to a clean, watertight shell at TSDF
parity.

% ─────────────────────────────────────────────────────────────────────────────
\subsection{Unreal Engine 5.8 Delivery}
\label{sec:tech:ue}

The game-engine delivery target is \emph{Unreal Engine~5.8}~\cite{epicgames2024_ue5},
Epic Games' production-grade real-time renderer.  Vitrine imports polygonal FBX
assets as Nanite-enabled static meshes with baked-texture materials.  The scene is
assembled via a Python control bridge over UE's Web Remote Control API
(\texttt{:30010}) and the experimental first-party UE MCP server (\texttt{:8000}).
An optional NanoGS~\cite{nanogs2026} splat layer provides a volumetric preview
alongside the polygonal assets, using Nanite-inspired LOD clustering to maintain
interactive frame rates.

The UE overlay runs inside a Docker container whose base image is
\texttt{nvidia/cuda:12.8.1-runtime-ubuntu22.04}, with the UE~5.8 Linux installed
build bind-mounted read-only from the host repository.  No Epic GitHub or
\texttt{ghcr.io} registry access is required.

\paragraph{USD (archival only).}
Universal Scene Description~\cite{pixar2019_usd} is used solely for archival export
and inter-tool exchange within the pipeline.  It is \emph{not} the UE deliverable:
UE's USD~Stage~Actor imports the live USDA and preserves Vitrine's \texttt{v2g:*}
custom metadata lineage tags via the \texttt{pxr} Python library, but LichtFeld's
native USD export emits a \texttt{ParticleField} prim that UE cannot import as
collision-ready geometry (ADR-019).

% ─────────────────────────────────────────────────────────────────────────────
\subsection{Summary of Component Choices}
\label{sec:tech:summary}

Table~\ref{tab:tech_summary} maps each pipeline stage to the primary technology,
the fallback, and the key decision rationale.

\begin{table}[h]
\centering
\caption{Summary of technology choices by pipeline stage.}
\label{tab:tech_summary}
\begin{tabular}{llll}
\toprule
\textbf{Stage} & \textbf{Primary} & \textbf{Fallback} & \textbf{Key driver} \\
\midrule
Feature matching   & ALIKED + LightGlue \cite{zhao2023_aliked,lindenberger2023_lightglue}   & SIFT + SuperGlue         & Speed + accuracy \\
SfM                & COLMAP \cite{schoenberger2016sfm}                                        & ---                      & Established baseline \\
3DGS training      & LichtFeld Studio \cite{lichtfeld2026}                                    & gsplat                   & MCP control surface \\
Recon enhancement  & ArtiFixer3D \cite{delutio2026_artifixer}                                 & Difix3D\texttt{+}        & Coverage-limited scenes \\
Scene mesh         & CoMe \cite{radl2026_come}                                                & TSDF / Open3D \cite{zhou2018_open3d}  & Confidence-guided accuracy \\
Object mesh        & MILo \cite{guedon2025_milo}                                              & SuGaR \cite{guedon2024_sugar}         & In-the-loop detail \\
Object hull gen    & TRELLIS.2 \cite{xiang2025_trellis2}                                      & Hunyuan3D-2.1 \cite{hunyuan2025_21}   & MIT licence + PBR \\
Segmentation       & SAM / SAM~2 \cite{kirillov2023_sam,ravi2024_sam2}                        & ---                      & Promptable, class-agnostic \\
Frame QA           & MUSIQ \cite{ke2021_musiq}                                                 & Q-Align \cite{wang2023_qalign}        & NR-IQA speed \\
UV atlas           & xatlas \cite{xatlas2018}                                                 & ---                      & No dependencies \\
Deblur training    & BAD-Gaussians \cite{zhao2024_badgaussians}                               & ---                      & Bundle-adjusted exposure \\
Game delivery      & UE~5.8 + FBX \cite{epicgames2024_ue5}                                    & NanoGS splat \cite{nanogs2026}        & Nanite + collision \\
\bottomrule
\end{tabular}
\end{table}
